% ID: SEL-D0074-OB-E03
% Tanda S074.
\textbf{Parte I. a)} Ninguno corresponde a $(A\cup B)^c=A^c\cap B^c$. Al menos uno corresponde a $A\cup B$. Que ocurra $B$ sin que ocurra $A$ corresponde a $B\cap A^c$.

\textbf{b)} Leemos $P(A/B)$ como $P(A\mid B)$. Por la definición de probabilidad condicionada, $P(A\cap B)=P(A\mid B)P(B)=0{,}3\cdot0{,}5=0{,}15$. Al sumar las probabilidades de $A$ y $B$, restamos la intersección que se ha contado dos veces:
\[
 P(A\cup B)=0{,}5+0{,}5-0{,}15=\boxed{0{,}85}.
\]

\textbf{Parte II.} La proporción muestral es $\hat p=140/200=0{,}70$. Con 140 respuestas positivas y 60 negativas, usamos la aproximación normal para la proporción, suponiendo que la muestra es aleatoria. Al $99\%$ queda $0{,}005$ en cada cola y $z_{0{,}995}\approx2{,}5758$. El intervalo es
\[
 0{,}70\pm2{,}5758\sqrt{\frac{0{,}70\cdot0{,}30}{200}}
 \approx\boxed{[0{,}6165;\ 0{,}7835]}.
\]
